% !TeX program = lualatex
% !TeX encoding = utf8
% !TeX root = ../Thermod.tex

%=========================================================
\Opensolutionfile{answer}[\currfilebase/\currfilebase-Answers]
\Writetofile{answer}{\protect\section*{\nameref*{\currfilebase}}}
\chapter{Співвідношення калібрування і співвідношення взаємності}\label{\currfilebase}
\makeatletter
\def\input@path{{\currfilebase/}}
\makeatother
%=========================================================

% --------------------------------------------------------
\section{1}
% --------------------------------------------------------
Викладемо ще раз набір аксіом, що є основою термодинаміки. Їх практичний зміст зводиться до наступних постулатів:
\begin{enumerate}
	\item [а)] Нульовий закон постулює існування функції стану --- температури:
	\begin{equation}\label{4.1}
		dT = \lpdr{T}{x}{F}dx + \lpdr{T}{F}{x}dF; ~~~~ \delta A = Fdx; ~~~~ \oint dT = 0
	\end{equation}
	\item [б)] Перший закон постулює існування функції стану --- внутрішньої енергії:
	\begin{equation}\label{4.2}
		dU = \delta Q - \delta A; ~~~~ \Delta U = \integ{1}{2}dU; ~~~~ \oint dU=0
	\end{equation}
	По-суті це закон збереження енергії.
	\item [в)] Другий закон постулює існування функції стану --- ентропії:
	\begin{equation}\label{4.3}
		dS = \dfrac{\delta Q}{T}; ~~~~ \Delta S = \integ{1}{2}dS; ~~~~ \oint dS = 0
	\end{equation}
\end{enumerate}
Другий закон доповнюється законом невід'ємності зміни ентропії в ізольованих системах та нерівністю Клаузіуса:
\begin{equation}\label{4.4}
	dS \geq 0; ~~~~ S_2-S_1 \geq \integ{1}{2}\dfrac{\delta Q}{T}; ~~~~ \oint \dfrac{\delta Q}{T} \leq 0
\end{equation}
Окрім цих постулатів існує ще третій, який не вводить ніяких функцій стану і до якого ми повернемося пізніше.
% --------------------------------------------------------
\section{2}
% --------------------------------------------------------
Об'єднуючи перший і другий закони, отримуємо основне рівняння термодинаміки:
\begin{equation}\label{4.5}
	dU = TdS - \sum_{i}F_idx_i,
\end{equation}
де $\delta A = \sum_{i}F_idx_i$ в загальному випадку. Якщо речовина ізотропна і однорідна, і виконувана нею робота супроводжується лише змінами об'єму, то:
\begin{equation}\label{4.5a}\tag{4.5a}
	dU = TdS - PdV
\end{equation}
Наслідком того, що функція $U$ є функцією стану є те, що приріст внутрішньої енергії $dU$ --- повний диференціал. Це в свою чергу означає, що змішані другі похідні дорівнюють одна одній:
\begin{equation}\label{4.6}
	\lpdr{T}{V}{S} = - \lpdr{P}{S}{V}
\end{equation}
Це співвідношення --- один із прикладів співвідношень взаємності, які є зрештою наслідком того, що $T$, $U$ i $S$ --- функції стану. Співвідношення (\ref{4.7}) може бути записане з допомогою \emph{якобіанів} або \emph{визначників Якобі}:
\begin{equation}\label{4.6a}\tag{4.6a}
	\jacobi{T,S}{V,S} = \jacobi{P,V}{V,S},
\end{equation}
За умови, $\jacobi{P,V}{V,S}\neq0$ отримуємо \emph{співвідношення калібрування}:
\begin{equation}\label{4.7}
	\jacobi{T,S}{P,V} =1
\end{equation}

Нерівність нулю (чи нескінченності) якобіана (\ref{4.7}) означає, що між парами змінних $(T,S)$ та $(P,V)$ існує взаємно однозначне співвідношення. Геометрична суть якобіана полягає в тому, що його модуль дає коефіцієнт зміни елементарної площі при переході від $(T,S)$ до $(P,V)$ площини. Звичайно, це справедливо для будь-яких пар незалежних змінних. Наприклад, елементарні площі в змінних $(x,y)$ і $(r,\varphi)$ пов'язані відношенням (див. рис.):
\begin{equation}\label{4.A}\tag{A}
	dxdy = rd\phi dr
\end{equation}
З іншого боку якобіан дорівнює:
\begin{equation}\label{4.B}\tag{B}
	\jacobi{x,y}{r\phi}=\lpdr{x}{r}{\phi}\lpdr{y}{\phi}{r}-\lpdr{y}{r}{\phi}\lpdr{x}{\phi}{r}=\cos\phi\cdot r\cos\phi+r\sin\phi\cdot\sin\phi=r
\end{equation}
Порівнюючи (\ref{4.A}) і (\ref{4.B}), знайдемо, що коефіцієнт пропорційності між площами $r$ --- дійсно рівний модулю якобіана. Така геометрична інтерпретація якобіана дозволяє розкрити глибинну суть співвідношення калібрування (\ref{4.7}). Площа, що описується замкнутим циклом, на індикаторній діаграмі $(P,V)$ дорівнює площі, що описується замкнутим циклом на діаграмі $(T,S)$. Це твердження є геометричною інтерпретацією (\ref{4.5a}), застосованому до замкнутих циклів\footnote{
	Властивості якобіанів другого порядку по незалежним змінним х,у:
	$$\mathfrak{J}(U,V) = \jacobi{U,V}{x,y}=
	\left\lvert
	\begin{matrix}
		\lpdr{U}{x}{y} & \lpdr{U}{y}{x}\\
		\lpdr{V}{x}{y} & \lpdr{V}{y}{x}
	\end{matrix}
	\right\rvert =
	\lpdr{U}{x}{y}\lpdr{V}{y}{x}-\lpdr{V}{x}{y}\lpdr{U}{y}{x}
	$$
	Із означення випливає:
	\begin{enumerate}
		\item $$\jacobi{U,y}{x,y}=\lpdr{U}{x}{y};$$
		\item $$\mathfrak{J}(U,V) = -\mathfrak{J}(V,U);$$
		\item $$\mathfrak{J}(kU,V)=\mathfrak{J}(U,kV)=k\mathfrak{J}(U,V);$$
		\item $$\jacobi{U,V}{kx,y}=\jacobi{U,V}{x,ky}=\dfrac{1}{k}\jacobi{U,V}{x,y};$$
		\item $$\jacobi{U,V}{x,y}=\jacobi{U,V}{x,y}\jacobi{r,S}{r,S}=\jacobi{U,V}{r,S}\jacobi{r,S}{x,y};$$
		\item $$\mathfrak{J}(U,V)\mathfrak{J}(z,y)+\mathfrak{J}(V,z)\mathfrak{J}(U,y)+\mathfrak{J}(z,U)\mathfrak{J}(V,y)=0~~(\text{\emph{тотожність Якобі}}).$$
	\end{enumerate}
	що дозволяє оперувати з ними як з дробами.

}.  Безумовно співвідношення калібрування (\ref{4.7}) може бути узагальнено і на інші термодинамічні змінні:
\begin{equation}\label{4.7a}\tag{4.7a}
	\jacobi{T,S}{F,x}=1
\end{equation}
Геометрична інтерпретація кількості тепла на діаграмі $(T,S)$ аналогічна для роботи на індикаторній діаграмі. Зважаючи на цю схожість, Гіббс запропонував називати ентропію термодинамічною координатою, а температуру --- термодинамічною силою. З практичної точки зору розгляд діаграми $(T,S)$ іноді значно зручніший, ніж $(P,V)$. Наприклад, незалежність К.К.Д. машини Карно від природи речовини проявляється в тому, що К.К.Д. на діаграмі $(T,S)$ є відношенням площ ABCD та ABC'D'  і, природно, від природи речовини не залежить.
% --------------------------------------------------------
\section{3}
% --------------------------------------------------------
Із загальних властивостей якобіана із (\ref{4.7}) випливає, що:
\begin{equation}\label{4.8}
	\jacobi{T,S}{P,V} = \jacobi{kT,k^{-1}S}{P,V}=1
\end{equation}
тобто існує довільність у виборі температурної і ентропійної шкал. Якщо температурну шкалу розтягнути в $k$ разів, то ентропійна шкала повинна бути стисненою в стільки ж. Із того, що незалежно від шкал кількість теплоти в термодинамічному процесі повинна бути однаковою, ця властивість стає практично очевидною:
\begin{equation}\label{4.9}
	dQ = T_1dS_1=T_2dS_2
\end{equation}
Із (\ref{4.9}) також випливає, що ентропія визначається з точністю до константи:
\begin{equation}\label{4.10}
	S_2=\dfrac{1}{k}S_1+S_0
\end{equation}
% --------------------------------------------------------
\section{4}
% --------------------------------------------------------
Тепер ми можемо узагальнити поняття термодинамічних коефіцієнтів, як величин:
\begin{equation}\label{4.11}
	\lpdr{\lambda}{\mu}{\nu}\text{, де } \lambda,~\mu,~\nu\text{ --- одна із величин } P,~V,~T,~S.
\end{equation}
Три з них пов’язані з ізотермічними і ще три з адіабатичними коефіцієнтами розширення $(\alpha_T,\alpha_S)$, тиску $(\beta_T,\beta_S)$ і стиснення $(\gamma_T,\gamma_S)$. Похідні $S$ по $T$ характеризують теплоємності:
\begin{equation}\label{4.12}
	C_V = T\lpdr{S}{T}{V}; ~~~~ C_P = T\lpdr{S}{T}{P}
\end{equation}
Похідні $S$ по $P$ вказують на необхідну кількість тепла, яке потрібно підвести до речовини, щоб залишились незмінними її об'єм чи температура під дією навантаження
\begin{equation}\label{4.13}
	\lpdr{S}{P}{V}= \dfrac{1}{T}\lpdr{Q}{P}{V}; ~~~~ \lpdr{S}{P}{T}=\dfrac{1}{T}\lpdr{Q}{P}{T}
\end{equation}
В свою чергу, похідні $S$ по $V$ визначають кількість теплоти, яку потрібно підвести до речовини, щоб тиск чи температура залишалися незмінними при її деформації
\begin{equation}\label{4.14}
	\lpdr{S}{V}{T} = \dfrac{1}{T}\lpdr{Q}{V}{T}; ~~~~
	\lpdr{S}{V}{P} = \dfrac{1}{T}\lpdr{Q}{V}{P}
\end{equation}
Складемо таблицю коефіцієнтів так, щоб перший рядок не мав коефіцієнтів $S$, другий --- $P$, третій --- $V$, четвертий --- $T$:
\begin{equation}\label{4.15}
	\begin{matrix}
		\lpdr{V}{P}{T} & \lpdr{P}{T}{V} & \lpdr{T}{V}{P} \\
		\lpdr{S}{T}{V} & \lpdr{T}{V}{S} & \lpdr{V}{S}{T} \\
		\lpdr{S}{T}{P} & \lpdr{T}{P}{S} & \lpdr{P}{S}{T} \\
		\lpdr{V}{P}{S} & \lpdr{P}{S}{V} & \lpdr{S}{V}{P} \\
	\end{matrix}
\end{equation}
Із 12 коефіцієнтів тільки три є незалежні. Дійсно, добуток коефіцієнтів в рядку дорівнює 1. Фактично, це наслідок того, що $T$ та $S$ --- функції стану і, отже, існують співвідношення типу $T = T(P,V) = T(S,V) = T(S,P)$ і $S = S(T,V) = S(T,P) = S(P,V)$. Переконаємося, що добуток двох будь-яких похідних, взятих при певній постійній величині, дорівнює третьому коефіцієнту:
\begin{equation}\label{4.16}
	\lpdr{V}{P}{S}\lpdr{P}{T}{S} = \jacobi{V,S}{P,S}\jacobi{P,S}{T,S}=\jacobi{V,S}{T,S}\jacobi{P,S}{P,S}=\lpdr{V}{T}{S};
\end{equation}
Співвідношення першого типу і типу (\ref{4.16}) зменшують кількість незалежних коефіцієнтів до 4, а співвідношенні калібрування залишає всього три незалежних коефіцієнти. Інші три коефіцієнти знаходять з експерименту або з використанням методів статистичної фізики.
% --------------------------------------------------------
\section{5}
% --------------------------------------------------------
Якщо говорити про експериментальну сторону питання, то легко виміряти наступні коефіцієнти:
\begin{equation}\label{4.17}
	\gamma_T = -\dfrac{1}{V}\lpdr{V}{P}{T}; ~~~~ \beta_V=\dfrac{1}{P}\lpdr{P}{T}{V}; ~~~~
	C_V = T\lpdr{S}{T}{V}
\end{equation}
Інші коефіцієнти можуть бути виражені через (\ref{4.17}) з використання техніки якобіанів\footnote{Необов'язково якобіанів, просто з ними зручніше.}. Для цього виключаємо ентропію із співвідношень з допомогою калібрування (\ref{4.7}), або ж відповідний якобіан виражається через теплоємності:
\begin{equation}\label{4.18}
	C_V = T\jacobi{S,V}{T,V}; ~~~~ C_P=T\jacobi{S,P}{T,P}
\end{equation}
З урахуванням (\ref{4.18}) і зроблених зауважень, отримаємо:
\begin{equation}\label{4.19}
	\lpdr{P}{S}{V}=\jacobi{P,V}{S,V}=\jacobi{P,V}{S,V}\jacobi{T,S}{P,V}=\jacobi{P,V}{P,V}\jacobi{T,S}{S,V}=-\lpdr{T}{V}{S};
\end{equation}
\begin{equation}\label{4.20}
	\lpdr{S}{V}{T}=\jacobi{S,T}{V,T}=\jacobi{S,T}{V,T}\jacobi{P,V}{T,S}=\jacobi{S,T}{T,S}\jacobi{P,V}{V,T}=\lpdr{P}{T}{V};
\end{equation}
\begin{equation}\label{4.21}
	\lpdr{S}{P}{T}=\jacobi{S,T}{P,T}\jacobi{P,V}{T,S}=\jacobi{S,T}{T,S}\jacobi{P,V}{P,T}=-\lpdr{V}{T}{P};
\end{equation}
\begin{equation}\label{4.22}
	\lpdr{T}{P}{S}=\jacobi{T,S}{P,S}=\jacobi{T,S}{P,S}\jacobi{P,V}{T,S}=\lpdr{V}{S}{P}.
\end{equation}
Співвідношення $(\ref{4.19})\div(\ref{4.22})$ отримали назву \enquote{\emph{співвідношення взаємності}} або \enquote{\emph{співвідношення Максвелла}}. Їх важливість полягає в тому, що вони пов'язують величини, які неможливо або складно виміряти безпосередньо в експерименті. Зауважимо, що у подібний спосіб коефіцієнт $\alpha$ можуть бути виражені через $\gamma$ і $\beta$:
\begin{equation}\label{4.23}
	\alpha=\gamma\beta P,
\end{equation}
а теплоємність $C_P$ пов'язана з $C_V$ співвідношенням (\ref{3.30}). Подібним чином можна знайти і інші коефіцієнти чи виразити коефіцієнти в $(\ref{4.19})\div(\ref{4.22})$ через інші величини. Наприклад:
\begin{equation}\label{4.24}
	\lpdr{T}{V}{S}=\jacobi{T,S}{V,S}\jacobi{P,V}{P,V}\jacobi{V,T}{V,T}=\jacobi{T,S}{P,V}\jacobi{P,V}{V,T}\jacobi{V,T}{V,S}=-\dfrac{T}{C_V}\lpdr{P}{T}{V}
\end{equation}
\begin{equation}\label{4.25}
	\lpdr{T}{P}{S}=\jacobi{T,S}{P,S}\jacobi{P,V}{P,V}\jacobi{P,T}{P,T}=\jacobi{T,S}{P,V}\jacobi{P,V}{P,T}\jacobi{P,T}{P,S}=-\dfrac{T}{C_P}\lpdr{V}{T}{P}
\end{equation}
% --------------------------------------------------------
\section{6}
% --------------------------------------------------------
Отримані формули можна експериментально перевірити. Особливу інтерес становлять перевірка співвідношень, при виведенні яких використовується співвідношення калібрування (\ref{4.7}), що є наслідком другого принципу. Один із подібних експериментів був проведений Джоулем і полягав у вивченні зміни температури адіабатично стискуваних рідин. Як випливає із (\ref{4.25}) при адіабатичному стисненні рідини на $dP$ її температура міняється на $dT$, що дорівнює:
\begin{equation}\label{4.26}
	dT = \dfrac{1}{C_P}T\lpdr{V}{T}{P}dP=\dfrac{\alpha VT_0}{C_P}dP.
\end{equation}
Всі величини в (\ref{4.26}) доступні експериментальній перевірці а, отже, (\ref{4.26}) є прекрасною ілюстрацією застосовності другого закону. Джоуль працював з риб'ячим жиром і водою і отримав доволі задовільне збіг експериментальних результатів із теоретичними. Так, в експериментах по стисненню води він встановив:
\begin{table}[!h]
	\centering
	\begin{tabular}{|c|c|c|c|}
		\hline
		$\Delta P$, бар & $T_0$, К & $\Delta T_{\text{експ.}}$ & $\Delta T_{\text{розрах.}}$ \\ \hline
		25,9    & 1,20  & -0,0083 & -0,0071   \\ \hline
		25,9    & 5,00  & 0,0044  & 0,0027    \\ \hline
		25,9    & 11,69 & 0,0205  & 0,0197    \\ \hline
	\end{tabular}
\end{table}
Показово, що експериментально було встановлено зменшення температури при адіабатичному стисканні рідини при $T<4\celsius$, що безумовно пов'язано з від'ємним значенням $\alpha$ при цих температурах. Без другого закону передбачити результати цих експериментів неможливо. Зрозуміло, що завдяки роботі, виконаній над рідиною, її внутрішня енергія зросте, проте як це вплине на температуру, \emph{апріорі} сказати неможливо, оскільки невідома залежність внутрішньої енергії від об'єму. Подібні досліди були проведені Хагою по розтягненню дротів. Для розрахунків скористаємося формулою:
\begin{equation}\label{4.27}
	dT = \lpdr{l}{T}{S}\dfrac{T_0}{C'_P}dF,
\end{equation}
де $C'_P$ --- теплоємність, віднесена до одиниці довжини, а $F$ --- сила, прикладена до дротини. В
дослідах зі сталлю і германським сріблом (мельхіором) було знайдено вражаючий збіг між
експериментальними результатами з теорією. Окрім цього, було знайдено механічний еквівалент теплоти,
що дорівнює 4,295 і 4,200 відповідно ($C_P$ вимірювалось в калоріях, а сили --- звичайно, в механічних
одиницях. Цікаво, що не особливо сильно розтягнута гума має від'ємний коефіцієнт теплового розширення.
Тоді можна стверджувати, що при адіабатичному розтягненні вона нагріється, що легко виявляється на
досліді простим дотиком губ.

% --------------------------------------------------------
\section{7}
% --------------------------------------------------------

Ефект зміни температури при адіабатичному стисненні (розширенні) речовини широко використовується для
отримання низьких температур і зрідження газів. Хоча шотландський фізик Джеймс Дьюар отримав рідкий
водень в 1898 р., отримання рідкого водню в значних кількостях вдалося налагодити тільки
Камерлінг-Оннесу. А в 1908 йому вдалося зрідити гелій --- один із так званих \enquote{вічних газів}.
Для цього використовувався ефект Джоуля-Томсона. Методом якобіанів можна довести, що коефіцієнт
диференціального ефекту Джоуля-Томсона становить:
\begin{equation}
	\lpdr{T}{P}{H}=\dfrac{V}{C_P}(\alpha T-1).
\end{equation}
і який для реальних газів $\Delta T\neq0$. а допомогою рідкого гелію Камерлінг-Оннесу вдалося досягти
ще нижчих температур: $1.38$ K 1909 року і $1.04$~K 1910-го. Це дозволило йому  провести низку
дослідженнь
властивостей речовин при таких низьких температурах. Своє найвідоміше відкриття Камерлінг-Оннес зробив
1911 року. Він встановив, що при низьких температурах електричний опір деяких металів повністю зникає,
зокрема у ртуті як видно на рис.~\ref{tikz:4.3Supercont}. Це явище Камерлінг-Оннес назвав
надпровідністю.

%---------------------------------------------------------
\begin{figure}[h!]\centering
    \begin{tikzpicture}
\begin{axis}[
    width=0.5\linewidth,
    height=0.75\linewidth,
    xlabel={Температура, K},
    ylabel={Опір, Ом},
    xmin=4.0, xmax=4.4,
    ymin=0, ymax=0.15,
    xtick={4.0,4.1,4.2,4.3,4.4},
    ytick={0.025,0.05,0.075,0.10,0.125,0.15},
    % Стиль под старину
    axis background/.style={fill=white},
    axis line style={thick},
    tick style={thick},
    grid=both,
    grid style={line width=.1pt, draw=gray!50, dotted},
    x tick label style={font=\small, /pgf/number format/fixed, /pgf/number format/precision=2},
    y tick label style={font=\small, /pgf/number format/fixed, /pgf/number format/precision=3},
    label style={font=\small},
    title style={font=\bfseries},
    every axis plot/.append style={thick},
    % Убираем современные элементы
    axis on top=true,
    tick align=outside,
]

% Данные
\addplot [
    only marks,
    color=black,
    mark=*,
    mark size=1.5pt
    ]
table {
4.372778827977316	0.12970922882427308
4.333459357277883	0.12629582806573958
4.232892249527411	0.11378002528445007
4.202646502835539	0.0011378002528445006
4.192060491493384	0.00018963337547408344
4.181474480151229	0.00018963337547408344
};

% Аппроксимация
\addplot[mark=none, domain=4.232:4.4, black] {0.1127 * x - 0.363};

% Пунктирная линия
\draw[dash pattern={on 1pt off 1pt}, thick] (4.23, 0.1137) -- (4.203,0.001138);

% Аннотация старого стиля
  \node[pin={[pin distance=0.5cm]95:{$10^{-5}$ Ом}}] at (4.192, 0.00018) {};

% Рамка вокруг графика (старый стиль)
\draw[thick] (current axis.south west) rectangle (current axis.north east);

\end{axis}
\end{tikzpicture}
\caption{Графік експерименту Камерлінг Оннеса, що ілюструє перехід ртуті в надпровідний стан}
\label{tikz:4.3Supercont}
\end{figure}
%---------------------------------------------------------


Але більш ефективним виявився метод зрідження гелію, запропонований Капіцею, в якому попередньо ізотермічно стиснутий газ адіабатично розширявся, при цьому в силу (\ref{4.24})  охолоджувався. Очевидно, що метод адіабатичного розширення можна застосовувати для будь-якого газу (в тому числі і ідеального). На сьогодні він широко використовується для отримання низьких температур і для зрідження газів.